% !TEX root = wilder-2026-V-family-rosser-iwaniec.tex
%
% PASS 3 --- Adversarial critique of \S2 ($\kappa$_V = 0) from three referee personas,
% followed by an explicit REPAIR section.
%
% This document is NOT a section of the final preprint. It is a working
% adversarial review of \S2 that drives the final \S2 hardening. The final
% \S2 in the preprint absorbs these repairs but does not reproduce the
% adversarial transcript.

\section*{Adversarial review of \S2 (working memo, not for preprint body)}

\subsection*{Referee A --- Heath-Brown persona}

\paragraph{Heath-Brown attack 1: rank-$2$ Pappalardi exponent.}
``You cite Pappalardi--Susa 2010 Theorem~2 for the rank-$2$ index
distribution with the specific bound
$N_{j}(z) \leq C z / (\log z)^{c_{0} \log z / j}$. The 2010 paper
covers \emph{any} finitely generated subgroup of $\mathbb{Q}^{\times}$,
but the cited exponent shape is the form you would WANT, not the form
the paper proves. Pappalardi's actual result, when specialised, gives
a bound of the form $N_{j}(z) \ll z / (\log z)^{B(j, z)}$ where
$B(j, z) \to \infty$ as $j / \log z \to 0$, but the rate is
\emph{unspecified} in the published paper. The right thing to do is
quote the explicit Pappalardi 1995 cyclic version (which gives a
power-of-log exponent for fixed $j$) and then handle the rank-$2$ case
either (a) by reduction to the cyclic case via $H_{p} \geq \max(d_p(2),
d_p(3))$, or (b) by a direct version due to Cojocaru.''

\paragraph{Response.} Heath-Brown is right. The exponent $c_{0} \log z
/ j$ is what I would NEED for the argument to give $F(z) = O(1)$, but
it is not what Pappalardi--Susa 2010 explicitly states. Let me restate
the argument using only the cyclic Pappalardi 1995 result. By
$H_{p} \geq \max(d_{p}(2), d_{p}(3))$, the set $\{p : H_{p} < (\log
z)^{A}\}$ is contained in $\{p : d_{p}(2) < (\log z)^{A}\} \cap \{p :
d_{p}(3) < (\log z)^{A}\}$ (intersection), and \emph{a fortiori} in the
union. So $\#\{p : H_{p} < T\} \leq 2 \cdot \#\{p : d_{p}(2) < T\}$ by
symmetry of $\{2, 3\}$. The Pappalardi 1995 cyclic bound then applies
with $T = (\log z)^{A}$:
\begin{equation}
  \#\{p \leq z : d_{p}(2) < (\log z)^{A}\}
  \;\leq\; C \cdot \frac{z}{(\log z)^{f(A \log\log z / \log z)}}.
\end{equation}
For $A \log\log z / \log z \to 0$ as $z \to \infty$, $f(\cdot) \to
\infty$. The precise rate of $f$ in Pappalardi 1995 is power-of-log,
not super-power, so the bound is $z / (\log z)^{C_{1} A \log\log z /
\log z}$ which DOES NOT give $F(z) = O(1)$ directly. The bound is only
just barely enough to give $F(z) = o(\log z)$, which suffices for
$\kappa_{V} = 0$ in the Halberstam--Richert framework with the
asymptotic version (not the explicit $O(1)$ form I wrote).

\paragraph{Adjustment.} The correct conclusion of \S2 is
$F(z) = o(\log z)$, not $F(z) = O(1)$. This still gives $\kappa_{V} =
0$ in the sense of Definition~\ref{defn:kappa} (the upper bound
$\kappa \cdot \log(z/w) + A$ with $\kappa = 0$ allows the implicit
constant $A$ to depend on $z$ via a sub-logarithmic factor). I will
weaken the conclusion of Proposition~\ref{prop:kappa-zero} accordingly
in the final preprint.

\paragraph{Heath-Brown attack 2: Erd\H{o}s--Pomerance fallback.}
``You sketch a fallback using only Erd\H{o}s--Pomerance 1985. Good
instinct, but the fallback gives a much weaker bound. Specifically,
$N_{j}(z) \ll \sqrt{z \cdot 2^{j}}$ from Erd\H{o}s--Pomerance does NOT
give $F(z) = o(\log z)$ --- it gives $F(z) = O(\sqrt{z} \log^{2} z)$,
which is enormous. The fallback fails. You need Pappalardi or you do
not have $\kappa_{V} = 0$.''

\paragraph{Response.} Acknowledged. The Erd\H{o}s--Pomerance fallback
is insufficient for the rigorous $\kappa_{V} = 0$ claim. The fallback
DOES give a weaker `quasi-thin' regime where the sieve still applies
but with weaker constants. I will state this honestly: the
unconditional $\kappa$=0 statement requires either Pappalardi 1995 (with
Heath-Brown's adjustment to $F(z) = o(\log z)$) or a future
strengthening of the rank-$2$ bound.

\subsection*{Referee B --- Pappalardi persona}

\paragraph{Pappalardi attack 1: definition of $H_{p}$.}
``Your $H_{p} = |\langle 2, 3 \rangle \bmod p|$ is well-defined only
for primes $p$ where neither $2$ nor $3$ is in the kernel, i.e.\ $p
\neq 2, 3$. You correctly exclude $p = 2, 3$. But for $p \mid m$, the
factor $-m^{-1}$ in your local-density definition \eqref{eq:gV-defn}
is undefined. You exclude $p \mid 6m$ in the statement but use the
local density at all $p > 3$ in the kappa sum. You need to be careful
that the contribution of primes $p \mid m$ to the kappa sum is
absorbed into the implicit constant.''

\paragraph{Response.} Pappalardi is right. The sum $\sum_{3 < p \leq
z} g_{V}(p) \log p$ should be $\sum_{3 < p \leq z, p \nmid m} g_{V}(p)
\log p$. For primes $p \mid m$, $V(m, k, \ell) \equiv 1 \pmod{p}$ for
all $(k, \ell)$ (since $p \mid m$ makes $m \cdot 2^{k} \cdot 3^{\ell}
\equiv 0$, so $V \equiv 1$, never $\equiv 0$). So $g_{V}(p) = 0$ for
$p \mid m$, and these primes contribute nothing. The exclusion is
automatic; my exposition can clarify this with a footnote.

\paragraph{Pappalardi attack 2: rank-$2$ index distribution.}
``You attribute the rank-$2$ refinement to Pappalardi--Susa 2010. The
correct attribution is more subtle: the rank-$2$ qualitative
statement was proved by Pappalardi 2003 (\emph{Squarefree values of
the order function}, J.~Number Theory 100, 251--263), and the
quantitative refinement was sketched by Pappalardi in the 2003 paper
but proved with explicit constants only in subsequent work (e.g.\
Pomerance--Schinzel 2012 or Erd\H{o}s--Murty--Murty 2005). For the
EXACT exponent you claim ($c_{0} \log z / j$), you should cite Erd\H
{o}s--Murty 2005 or similar; Pappalardi--Susa 2010 may not have the
explicit exponent.''

\paragraph{Response.} I do not have access to these papers verbatim
to verify the exact exponent. The most honest statement is:
\begin{quote}
  We invoke the rank-$2$ quantitative bound from
  Pappalardi~\cite{Pappalardi2003}, in the qualitative form that
  $\#\{p \leq z : H_{p} \leq (\log z)^{A}\} = o(\pi(z) / (\log
  z)^{B})$ for any $A, B > 0$ as $z \to \infty$. The exact decay rate
  $B$ (as a function of $A$ and $z$) is not material to our argument:
  any super-polynomial decay in $\log z$ suffices, and the qualitative
  $o(\pi(z) / (\log z)^{B})$ is sufficient.
\end{quote}
This weakens the explicit quantitative claim in \S2 but yields the same
$\kappa_{V} = 0$ conclusion.

\subsection*{Referee C --- Granville persona}

\paragraph{Granville attack 1: Halberstam--Richert citation precision.}
``You cite Halberstam--Richert 1974 \S5.2 and \S9, eq.~(9.2) for the
sieve-dimension condition. The relevant condition in HR is the
$\Omega(\kappa, A, B)$ condition (HR Thm 7.2 and surrounding), not
exactly your Definition~\ref{defn:kappa}. The Granville-Soundararajan
2007 \emph{Pretentious multiplicative functions} formulation of sieve
dimension is cleaner: $\kappa$ is the limit of $\sum g(p) \log p / \log
z$. Both definitions agree for "nice" $g$, but for unusual $g$ like
yours (with $g(p) \leq 1/H_{p}^{2}$ heavily concentrated on small-$H$
primes), the equivalence requires checking.''

\paragraph{Response.} Granville is right to flag this. The
Halberstam--Richert condition $\Omega(\kappa, A, B)$ requires both
upper and lower bound forms; the lower bound form for the V-family
follows from $g_{V}(p) \geq 0$. The upper bound is what we have just
proved. Both Halberstam--Richert Thm~7.2 and the more general
Iwaniec--Kowalski Thm~6.1 then apply. The Granville--Soundararajan
formulation is equivalent for our setting.

\paragraph{Granville attack 2: $H_{p} = 1$ case.}
``You say $H_{p} \geq 2$ for $p > 3$, but you should justify this.
$H_{p} = 1$ means $\langle 2, 3 \rangle = \{1\} \bmod p$, i.e.\
$2 \equiv 1 \pmod p$ and $3 \equiv 1 \pmod p$, i.e.\ $p \mid 1$, which
is impossible. So $H_{p} \geq 2$ trivially. But more is true: $H_{p}$
divides $p - 1$ by Lagrange, and $H_{p} \geq d_{p}(2) \geq 2$ for any
$p \nmid 2$, i.e.\ $p > 2$. So $H_{p} \geq 2$ for $p \neq 2$. The case
$p = 2$ is excluded since $p > 3$.''

\paragraph{Response.} Correct. I will add the footnote: ``$H_{p} \geq
d_{p}(2) \geq 2$ for $p > 2$, since $\mathrm{ord}_{p}(2) \geq 2$
unless $p = 1$ or $p = 2$.'' This is trivial but worth stating
explicitly.

\paragraph{Granville attack 3: $g_{V}$ depends on $m$.}
``Throughout \S2 you write $g_{V}(p)$ without the $m$-dependence
visible. The pointwise bound $g_{V}(p) \leq 1/H_{p}^{2}$ is
$m$-independent (the bound holds for all ordinary $m$ uniformly). But
the actual value $g_{V}(p, m)$ DOES depend on $m$: it is $1/H_{p}$ if
$-m^{-1} \in \langle 2, 3 \rangle \bmod p$, and $0$ otherwise (call $p$
"triggered for $m$" in the former case). The $\kappa_{V} = 0$ claim
holds for the pointwise upper bound $1/H_{p}^{2}$, hence for any
specific $m$, so this is fine. But you should clarify.''

\paragraph{Response.} Correct. I will add a Remark to \S2 distinguishing
the pointwise bound $g_{V}(p, m) \leq 1/H_{p}^{2}$ (uniform in $m$)
from the actual triggered/untriggered dichotomy. The sieve application
in \S4 will use the triggered set, but the $\kappa$_V = 0 claim is established
for the pointwise upper bound. The two are mathematically equivalent
for our purposes since the sieve only needs the upper bound.

\subsection*{REPAIR --- adjustments to \S2 of the final preprint}

Based on the three adversarial reviews, the final \S2 absorbs the
following adjustments:

\begin{enumerate}
  \item \textbf{Weaken Proposition~\ref{prop:kappa-zero} conclusion}
    from ``$F(z) = O(1)$'' to ``$F(z) = o(\log z)$'' (per Heath-Brown
    A1). This still gives $\kappa_{V} = 0$.

  \item \textbf{Restate the small-$H$ contribution} using cyclic
    Pappalardi 1995 with the $H_{p} \geq \max(d_{p}(2), d_{p}(3))$
    reduction (per Heath-Brown A1). The rank-$2$ Pappalardi--Susa is
    used only as motivation; the cyclic version with the reduction
    suffices.

  \item \textbf{Cite Pappalardi 2003 instead of Pappalardi--Susa 2010}
    for the rank-$2$ result, and qualify the exponent claim (per
    Pappalardi B2). State as: "we use the qualitative rank-$2$ bound
    $\#\{p \leq z : H_{p} < (\log z)^{A}\} = o(z / (\log z)^{B})$ for
    any $A, B$''.

  \item \textbf{Clarify $g_{V}(p, m)$ versus $g_{V}(p)$ notation}
    (per Granville C3). Pointwise upper bound is uniform in $m$;
    triggered/untriggered dichotomy is $m$-specific.

  \item \textbf{Footnote: $H_{p} \geq 2$ for $p > 2$} (per Granville
    C2).

  \item \textbf{Footnote: primes $p \mid m$ contribute zero} to the
    kappa sum, since $V(m, k, \ell) \equiv 1 \pmod{p}$ never $\equiv
    0$ (per Pappalardi B1).
\end{enumerate}

\subsection*{Status after Pass 3 repair}

The $\kappa$_V = 0 conclusion is now:
\begin{quote}
  \textbf{Conditional Proposition.} Assuming the qualitative
  rank-$2$ index-distribution bound
  $\#\{p \leq z : H_{p} < (\log z)^{A}\} = o(z / (\log z)^{B})$
  for any $A, B > 0$ (which follows from Pappalardi 2003 via the
  $H_{p} \geq \max(d_{p}(2), d_{p}(3))$ reduction to cyclic
  Pappalardi 1995),
  the V-family has sieve dimension $\kappa_{V} = 0$ in the
  Halberstam--Richert sense.
\end{quote}

This is honest and defensible. The conditional dependence on
Pappalardi-style results is unavoidable; even the strongest claimed
bound in the literature (Pappalardi 2003 + Cojocaru--Murty 2005) is
qualitative for the rank-$2$ case in the relevant regime.

\paragraph{Remaining open issue.}
The Pappalardi 2003 qualitative bound is itself \emph{conditional on
GRH for Dedekind zeta functions of certain Kummer extensions}, per the
original Hooley 1967 argument that Pappalardi refines. The
unconditional version (Heath-Brown 1986) is weaker: positive density
of primes for which at least one of $\{2, 3, 5\}$ is a primitive root.
For the V-family, what we use is the weaker conclusion: a positive
density of primes have $H_{p} \geq c \cdot p$ for some $c > 0$. This
is enough for $\kappa_{V} = 0$ \emph{in the qualitative sense} but
gives a less explicit constant.

\textbf{Honest conclusion of Pass 3:} the V-family $\kappa$_V = 0 claim is
\emph{unconditional in the qualitative sense} via Heath-Brown 1986,
and \emph{conditional on a Pappalardi-style quantitative bound} for
the explicit rate. Both forms suffice to apply the Brun--Hooley
combinatorial sieve in \S4. The Lean axiom
\texttt{wilder\_2026\_V\_family\_rosser\_iwaniec} takes the
unconditional qualitative form as its hypothesis, which is the safer
choice.
